El presente proyecto tuvo como propósito mejorar y automatizar un brazo robótico cartesiano existente para su aplicación en la línea de empaque del Laboratorio de Diseño de Procesos de la Universidad del Valle de Guatemala. El sistema presentaba limitaciones mecánicas en los ejes X, Y y Z, tales como fricción, desalineaciones y movimientos irregulares, que afectaban su funcionamiento y no se corrigieron mediante lubricación.

Para atender estas deficiencias se realizó una reingeniería parcial de la estructura y de los sistemas de traslación mediante la incorporación de ejes guía de acero y rodamientos lineales. Además, se integraron finales de carrera, un encoder en la banda transportadora, una cámara HUSKYLENS 2 con un modelo de segmentación y una arquitectura de control distribuida entre una Portenta H7 con Machine Control y una ESP32. El sistema incorporó un modo manual con mando inalámbrico, una rutina de referenciación automática, una interfaz local en pantalla OLED y un modo automático capaz de recoger piezas sobre la banda en movimiento y trasladarlas a un punto de entrega.

El trabajo se organizó en seis fases: evaluación del sistema original, rediseño mecánico, implementación mecánica, accionamiento electrónico y operación manual, integración de visión y control automático, y validación experimental. Durante la integración, la comunicación entre los microcontroladores se migró de I2C a UART sobre RS485 debido a la corrupción persistente de paquetes.

La validación consistió en 120 ciclos automáticos de recogida y entrega con dos piezas, tres orientaciones y dos velocidades de banda. El 61,7~\% de los ciclos fue exitoso, con mejor desempeño a menor velocidad y con las piezas orientadas a lo largo de la banda. La referenciación mostró una dispersión del recorrido medido de 1,3~mm en X y 0,5~mm en Y. La exactitud y la repetibilidad de posicionamiento del efector final no se midieron, por lo que el cuarto objetivo se cumplió de forma parcial.
